\documentclass[11pt]{letter}
\usepackage[margin=1in]{geometry}
\usepackage{times}
\usepackage{xcolor}

\signature{Jeff Gore (Corresponding Author), Jinyeop Song, and Jiliang Hu \\ Department of Physics \\ Massachusetts Institute of Technology \\ Cambridge, MA, USA \\ jeffgore@gmail.com}

\address{Department of Physics \\ Massachusetts Institute of Technology \\ Cambridge, MA 02139, USA}

\begin{document}

\begin{letter}{}

\opening{Dear Editor,}

We are pleased to submit our manuscript, ``Interspecies Interactions Drive Community-Level Selection in Microbial Coalescence,'' for consideration in Nature Ecology \& Evolution.

A fundamental unresolved question in ecology is whether communities function as cohesive units of selection (Clements, 1916) or loose collections of independent species (Gleason, 1926). Community coalescence provides a natural test, yet empirical studies show conflicting results: some demonstrate community-level selection, while others observe independent species sorting. A theoretical framework explaining when and why each regime arises is missing.

In this work, we combine ecological modeling with over 300 coalescence experiments among randomly assembled bacterial communities to demonstrate that interspecies interaction strength is the governing parameter that determines the level of selection. We show that weak interactions lead to independent species sorting, whereas moderate-to-strong interactions drive correlated species persistence and community-level selection. We validate this experimentally by manipulating nutrient concentrations to tune interaction strength, shifting the system between these two regimes in alignment with our theoretical predictions. Furthermore, we identify two distinct mechanisms driving this cohesion, one emergent and one top-down, and demonstrate that these interaction-dependent patterns generalize to complex communities derived from diverse natural environments.

Crucially, our framework reconciles conflicting observations in the recent literature by identifying interaction strength as the ``tuning knob'' for selection. This provides a unified framework for community assembly that explains why different environments favor different selection regimes. While demonstrated in bacterial microcosms, the underlying framework is not constrained to microbial ecology and may extend to plant, animal, and other communities. We believe this manuscript will be of significant interest to the readership of Nature Ecology \& Evolution.

We have not had any prior discussions with a Nature Ecology \& Evolution editor about this work. This manuscript has not been published and is not under consideration for publication elsewhere, and no related manuscripts are under consideration or in press elsewhere. We have no conflicts of interest to disclose.

We would like to suggest the following reviewers for this manuscript:
\begin{itemize}
\item \textcolor{red}{[Reviewer Name], [Institution] --- [brief reason for suggestion]}
\item \textcolor{red}{[Reviewer Name], [Institution] --- [brief reason for suggestion]}
\item \textcolor{red}{[Reviewer Name], [Institution] --- [brief reason for suggestion]}
\end{itemize}

Thank you for considering our manuscript.

\closing{Sincerely,}

\end{letter}

\newpage
\noindent{\Large\textcolor{red}{\textbf{Notes for Jeff (to be removed before submission):}}}

\begin{enumerate}
\color{red}
\item \textbf{Related manuscripts disclosure.} The journal requires us to disclose any related manuscripts that any of the authors have currently under review or in press elsewhere. ``Related'' means other papers by us that cover overlapping topics, use the same data, or describe methods used in this study (e.g., a companion paper, a methods paper, or a review on coalescence). Do we have any such manuscripts to disclose? If yes, we need to briefly describe them in the cover letter.

\item \textbf{Suggested reviewers.} We need to fill in the suggested reviewers section (names, institutions, and a brief reason for each suggestion). We can also optionally list anyone we would like excluded from reviewing, with a reason. Please provide 3--5 names.
\end{enumerate}

\end{document}
